\section{Formalization}\label{sec:12}
\textbf{What the kernel certifies.} Apart from the literature-to-$\mathrm{TCER}$ adapter in \cref{lemma:10.8}, every result stated above is a Lean 4 declaration, with the external inputs retained as the explicit hypotheses specified below; the declarations are in the accompanying artifact, archived at ⟨DOI⟩. For declarations cited in the paper-claim registry, \texttt{scripts/PaperClaimTypeReport.\allowbreak{}lean} records elaborated types and axiom footprints; \texttt{scripts/FinalAxiomReport.\allowbreak{}lean} is a separate report over its own listed roots. Every recorded footprint is contained in the standard set \texttt{[propext, Classical.\allowbreak{}choice, Quot.\allowbreak{}sound]}; contained, not equal, since the roots do not all use all three. The forbidden tokens \texttt{sorryAx}, \texttt{native\_\allowbreak{}decide}, \texttt{ofReduceBool}, \texttt{trustCompiler} do not occur anywhere in the development. The toolchain is \texttt{leanprover/lean4:v4.\allowbreak{}31.\allowbreak{}0-rc1} with the mathlib revision pinned in the lakefile; the Lean tree distributed with the paper is copied from a clean-room build whose source manifest is required to match hash for hash, and the artifact ships the build script and the manifests. The counts — roots, their division among subsets of the standard set, and the external-input-free roots — are in \texttt{release/TRUST\_\allowbreak{}BOUNDARY\_\allowbreak{}APAL.\allowbreak{}md} of the artifact, where a tool recomputes them from the certificates and fails the build if the document and the development disagree.

\textbf{External inputs are hypotheses.} The two classical inputs (\cref{definition:1.5,definition:1.6}) and the analytic input $\mathrm{TCER}$ of \cref{sec:10} are \texttt{def … : Prop}, and a result uses one only by taking it as an argument. $\mathrm{TCER}$ is not formalized: \cref{lemma:10.8} derives it from \cite{KK26} in ordinary mathematics. \texttt{\#print axioms} reports transitive \emph{axiom} dependencies and is blind to hypotheses, so a conditional and an unconditional theorem can print the same footprint; which results are unconditional is therefore read from a second report, \texttt{scripts/HypothesisBinderLedger.\allowbreak{}lean}, which inspects the hypothesis binders of each root's type and reports the roots that carry no external-input binder, explicit, implicit, instance or wrapped, with seven planted-hypothesis controls. A few load-bearing declarations — notably the conversion of a structured family into a certificate in \cref{sec:8} — are definitions producing data rather than theorems and so are not roots of the report; their footprints are inherited by every theorem that uses them, and those theorems are reported. The layers are:

\begin{longtable}{@{}>{\raggedright\arraybackslash}p{0.279\linewidth}>{\raggedright\arraybackslash}p{0.388\linewidth}>{\raggedright\arraybackslash}p{0.232\linewidth}@{}}
\toprule
layer & results & input used \\
\midrule\endhead
unconditional & \cref{theorem:2.1,theorem:2.2,theorem:3.2,theorem:3.3,theorem:3.4,theorem:4.1,theorem:5.1,theorem:6.1,theorem:6.2}, \crefrange{theorem:8.3}{theorem:8.5}, \cref{theorem:9.2,theorem:9.3} & none beyond mathlib \\
conditional on \cref{definition:1.5} & SE₂ ⊆ REC and its strictness at positive arity (\cref{theorem:6.3} and \cref{corollary:6.6}) & code-uniform quadratic solvability \\
conditional on \cref{definition:1.6} & RE = PP₂ = SE₄ and the function phase diagram of \cref{sec:7} & DPRM over ℕ \\
conditional on \cref{definition:1.5} and \cref{definition:1.6} & the complete relation phase diagram (\cref{theorem:6.8}) & quadratic solvability and DPRM over ℕ \\
unconditional, \cref{sec:10} & \cref{proposition:10.1,proposition:10.2,lemma:10.3,lemma:10.4,lemma:10.5,proposition:10.6,proposition:10.11,proposition:10.12,lemma:10.13,proposition:10.16,proposition:10.18}; the reverse directions of \cref{proposition:10.9,proposition:10.14,proposition:10.17} & none beyond mathlib \\
conditional on $\mathrm{TCER}$ & \cref{proposition:10.7,proposition:10.10,proposition:10.15}; the forward directions of \cref{proposition:10.9,proposition:10.14,proposition:10.17} & TCER (\cref{lemma:10.8}) \\
\bottomrule
\end{longtable}
$\mathrm{Fun}(SE_{2}) ⊆ \mathrm{TotComp}$ belongs to the first row and not the second: it follows from totality together with a Diophantine graph, and uses neither classical input nor the height bound.

\textbf{Statement fidelity.} A formal proof certifies that a formal statement follows from the axioms; it does not certify that the statement says what the prose says. Three mechanisms guard that gap, and none of them establishes the correspondence by itself. The principal formal statements are registered in a frozen-statement layer: each registered declaration is frozen in a separate file and diffed against the production tree as declaration text with runs of whitespace normalised (\texttt{scripts/statement\_\allowbreak{}diff.\allowbreak{}py}, with a negative control that plants a mutation and a missing declaration and requires both to be caught); definition bodies are covered separately by an elaborated-value hash certificate. The artifact's reconciliation report identifies the exported roots outside that registry. Second, every numbered theorem, proposition and corollary of this paper that names a Lean declaration has a row in a paper-claim registry (\texttt{submission/native/PAPER\_\allowbreak{}CLAIMS.\allowbreak{}tsv}) recording the Lean roots it cites, the quantifier pattern, the degree bound, the witness domain and the external hypotheses, together with a hash of the paper statement and the SHA-256 of each cited root's elaborated Lean type (\texttt{scripts/PaperClaimTypeReport.\allowbreak{}lean}). The registry detects unreviewed changes to registered paper statements and to Lean root types: editing a statement, substituting an unrelated root, or changing a root's type fails the build, and the artifact ships planted negative controls, all of which must be caught. The semantic correspondence between a mathematical statement and a Lean type is recorded one row per cited root (\texttt{PAPER\_\allowbreak{}CLAIMS\_\allowbreak{}REVIEW.\allowbreak{}tsv}) as a model-assisted, author-approved correspondence registry: each row was read by a language model against the fully elaborated Lean type and its disposition was approved by the author; it is not established by text hashing alone, and the correspondence remains an author-approved mapping supported by the type and axiom reports, not a machine-certified one. Third, additional validation includes the concrete coefficient examples formalized in \texttt{SemanticQuadraticTwoWitness/Fidelity2W.\allowbreak{}lean} and \texttt{SemanticQuadraticSingular/Fidelity3W.\allowbreak{}lean}, and the validation record of \cref{sec:C} (\texttt{certificates/APPENDIX\_\allowbreak{}C\_\allowbreak{}STATEMENT\_\allowbreak{}FIDELITY.\allowbreak{}log}), which reports a comparison of the monomial-count tables with an independent Python enumeration. These checks are supplementary and do not establish semantic correspondence between the paper and the Lean types; that correspondence remains subject to author review, recorded in the artifact (\texttt{release/AUTHOR\_\allowbreak{}SIGNOFF\_\allowbreak{}SHEET.\allowbreak{}md}).

\textbf{Names cited here.} Every declaration named in this paper exists in the development. The principal results and their roots are: the finite-or-structured theorem for single-quadratic fibres (\cref{theorem:8.5}, \texttt{uniformQuadratic\allowbreak{}FiberStructureZZ\_\allowbreak{}of\_\allowbreak{}quadratic}; its per-fibre corollary \texttt{singleQuadratic\_\allowbreak{}fiber\_\allowbreak{}bounded\_\allowbreak{}or\_\allowbreak{}structurallyInfinite} is the weaker ∀a∃B form of \cref{definition:8.4} and is not the theorem); the arithmetic projection dichotomy, bounded or infinite (\cref{theorem:9.2}, \texttt{quadraticProjectionDichotomyZ\_\allowbreak{}all}) and its data-carrying form (\cref{theorem:9.3}, \texttt{quadraticProjectionDichotomyZ\_\allowbreak{}all\_\allowbreak{}data}); the finite-fibre height bound (\cref{theorem:2.1}, \texttt{singleQuadratic\_\allowbreak{}finiteFiber\_\allowbreak{}height\_\allowbreak{}vector}); definable-function growth (\cref{theorem:3.2}, \texttt{singleQuadratic\_\allowbreak{}function\_\allowbreak{}quadratic\_\allowbreak{}growth}); the polynomial-graph classification (\cref{theorem:5.1}, \texttt{polynomial\_\allowbreak{}graph\_\allowbreak{}singleQuadratic\_\allowbreak{}iff}); $\mathrm{SE}_{2} ⊆ \mathrm{REC}$ under the Grunewald–Segal principle (\cref{theorem:6.3}, \texttt{singleEqFragmentZZ\_\allowbreak{}two\_\allowbreak{}subset\_\allowbreak{}recFragmentZZ}); the relation phase diagram at every positive arity under the two literature hypotheses (\cref{theorem:6.8}, \texttt{relation\_\allowbreak{}phase\_\allowbreak{}diagram\_\allowbreak{}all\_\allowbreak{}positive\_\allowbreak{}arity\_\allowbreak{}of\_\allowbreak{}classical}); and the function phase diagram under $\mathrm{DPRM}_{NN}$ (\cref{theorem:7.1}, \texttt{function\_\allowbreak{}phase\_\allowbreak{}diagram\_\allowbreak{}by\_\allowbreak{}arity\_\allowbreak{}of\_\allowbreak{}DPRM\_\allowbreak{}NN}); the exact roots of every numbered result are the rows of the paper-claim registry, and the full root map is in the artifact. The multi-parameter universal-pair question of \cref{sec:C} is an application of the same development and is not a headline result. Some names are cited as signposts into the sources and are covered by neither the frozen set nor the axiom report under those names; no numbered result of this paper rests on such a name. The census of cited names, and the reconciliation of the axiom report with the frozen set, are in the trust document, recomputed at every build.

Lean is an additional verification of the proofs given in this paper, not a substitute for them: every numbered result is proved in the text or in the appendices.

